% Generated draft: requires mathematical and source review.
\documentclass[11pt]{article}
\usepackage[margin=25mm]{geometry}
\usepackage{fontspec}
\setmainfont{Latin Modern Roman}
\usepackage{amsmath,amssymb,mathtools,unicode-math}
\setmathfont{Latin Modern Math}
\usepackage{enumitem,needspace,xurl,hyperref}
\setlength{\parindent}{0pt}
\setlength{\parskip}{5pt}
\setlength{\emergencystretch}{3em}
\newcommand{\N}{\mathbb{N}}
\newcommand{\Z}{\mathbb{Z}}
\newcommand{\Q}{\mathbb{Q}}
\newcommand{\R}{\mathbb{R}}
\newcommand{\C}{\mathbb{C}}
\newcommand{\F}{\mathbb{F}}
\newcommand{\A}{\mathbb{A}}
\newcommand{\PP}{\mathbb{P}}
\newcommand{\E}{\mathbb{E}}
\newcommand{\eps}{\varepsilon}
\newcommand{\dd}{\,\mathrm{d}}
\begin{document}
\section*{UnsolvedMath ID: 30003750}
\textbf{Source problem code:} catalogue entry
\subsection{Definitions and mathematical statement}
Let \(\gamma\) denote the dimension-to-sample ratio in the supplied record, with \(0\leq\gamma<1\). Write \(\mathcal P(\R)\) for the Borel probability measures on \(\R\), and let \(\mathcal A_\gamma\subseteq\mathcal P(\R)\) denote the intended class of attainable limiting empirical distributions of projections of structureless Gaussian data. The source record does not define this class precisely.

For \(q>0\), \(N(0,q)\) denotes the centered Gaussian probability law with variance \(q\). The supplied conjecture asks whether every \(\mu\in\mathcal A_\gamma\) admits a decomposition of the form
\[
\mu=(1-\eta)N(0,q)+\eta\nu,
\qquad q<(1+\sqrt\gamma)^2,
\qquad 0\leq\eta\leq1,
\qquad \nu\in\mathcal P(\R),
\]
where \(\nu\) is described as a non-Gaussian component when it has positive weight, and \(\eta\) is described as possibly small. Equality here means equality of probability measures: for every Borel set \(B\subseteq\R\),
\[
\mu(B)=(1-\eta)N(0,q)(B)+\eta\nu(B).
\]
The absent component corresponds to \(\eta=0\).

\textbf{Unresolved qualifications.} This displays the intended decomposition, but the supplied record does not specify a complete quantified conjecture. It gives no numerical meaning to ``small,'' no requirement excluding \(\eta=1\), and no rule governing dependence of the component weight on \(\gamma\) or \(\mu\). Allowing \(\eta=1\) makes the decomposition automatic for any non-Gaussian \(\mu\), by taking \(\nu=\mu\). It is also unclear whether variance zero is permitted; if it is, \(N(0,0)\) would mean the point mass at zero. The data normalization, asymptotic regime, admissible projection directions, dependence of directions on the observations, and mode of convergence defining \(\mathcal A_\gamma\) must be specified before the conjecture has a unique mathematical meaning. The strict upper bound on \(q\) is retained exactly as supplied.
\subsection{Short English statement}
For Gaussian data with dimension-to-sample ratio \(\gamma<1\), the supplied conjecture proposes that every attainable limiting projection distribution consists of a centered Gaussian component of variance below \((1+\sqrt\gamma)^2\), together with a possibly small non-Gaussian component. The attainable class and the permitted component weights are unspecified, so this remains a statement-only draft requiring mathematical review.
\subsection{Sources}
\begin{itemize}
\item Local registry: \texttt{data/unsolvedmath/problems.json}, with \texttt{catalog/problems.json} as the supplemental registry specified by the project; permanent UnsolvedMath ID \(30003750\), title ``Attainable Limits in Gaussian Projection Pursuit,'' exact source problem code \texttt{catalogue entry}. These are registry references, not independently inspected files.
\item Supplied report citation: \emph{Statistical Inference for Structured High-dimensional Models} (2019), Oberwolfach Reports, pp.~591--630, DOI \texttt{10.4171/owr/2018/12}. \url{https://doi.org/10.4171/owr/2018/12}. The registry gives proposed year 2018 and names no individual proposer; attribution is to the supplied workshop-report provenance.
\item Supplied literature source: \emph{Projection pursuit in high dimensions}. \url{https://doi.org/10.1073/pnas.1801177115}.
\item Supplied literature source: \emph{Projection pursuit perspective}. \url{https://doi.org/10.1111/insr.12517}.
\item Literature assessment dated \texttt{2026-08-21}: the supplied curation labels the problem open and reports Gaussian and near-Gaussian results, but no complete proof or disproof of the proposed characterization for all distributions in the subcritical regime. This is supplied dated curation, not an independently verified literature assessment. The background expressly requests independent verification of time-sensitive status and cited claims.
\item The supplied statement status is \texttt{corrected\_verified}; its accompanying explanation says mathematical content was preserved while notation, typography, and sentence structure were repaired for publication. This supplied editorial label does not resolve the formulation issues identified here or establish mathematical correctness.
\end{itemize}
\end{document}
